\section{Noncommutative geometry and physical interpretation}
\label{sec:ncg}
\subsection{Covariance rather than exact order}
For a family of configurations $\omega\in\Omega$, a spatial action
$\tau_a$ is represented by
\begin{equation}
 U_a H_\omega U_a^{-1}=H_{\tau_a\omega}.
 \label{eq:covariance}
\end{equation}
This covariance is weaker than $[U_a,H_\omega]=0$: a displaced thermal
configuration need not have the symmetries of its probability measure.
Bellissard's formulation replaces the Brillouin-zone description by
an algebra of such covariant observables with a trace per physical
volume \cite{Bellissard1994}. The operator-algebraic viewpoint therefore
remains meaningful when a single Bloch momentum is not a good quantum
number.

For sufficiently local observables, the trace per volume has the
schematic realization
\begin{equation}
 \mathcal T(A)=
 \lim_{V\rightarrow\infty}\frac{1}{V}
 \Tr(\chi_V A_\omega\chi_V),
 \label{eq:trace-volume}
\end{equation}
with the appropriate orbital trace and a physical-volume normalization.
Ergodicity supplies the conditions under which a typical large
configuration and an ensemble average agree. It does not guarantee
that a particular small simulation cell has self-averaged.

The first-principles construction must preserve covariance while
passing from continuum KS operators to AO matrices
\cite{KuehneTransport2020}. The metric is part of that transformation.
In a finite AO space, traces of projected covariant matrices use
$\Tr(S^{-1}A)$; traces of products use the intervening metric factors.
For local traces the spatial cutoff itself is an operator, not merely
an arbitrary subset of coefficient indices.

\subsection{Derivations, Fermi projections, and Chern pairings}
Choose spatial derivations
$\nabla_a A=-\ii[X_a,A]$ in an orthonormal representation.
For the Fermi projection $P=\mathbf1_{(-\infty,E_F)}(H)$, a
two-dimensional Chern pairing can be written, with a fixed orientation,
as
\begin{equation}
 C=2\pi\ii\,\mathcal T\!\left(P[\nabla_1P,\nabla_2P]\right)
   =-2\pi\ii\,\mathcal T\!\left(P[[X_1,P],[X_2,P]]\right).
 \label{eq:nc-chern}
\end{equation}
The trace normalization and oriented coordinates are chosen so that
the periodic limit agrees with the Berry-curvature convention used
below. Reversing that orientation reverses $C$.
Index theory relates such pairings to stable integers under the
required locality and gap or localization assumptions
\cite{ProdanSchulzBaldes2016}. Equation~\eqref{eq:nc-chern} is a theoretical
connection, not a claim that the current implementation evaluates a
general thermodynamic-limit noncommutative Chern formula.

The metric also enters a finite real-space evaluation. Let
$C_{\mathrm{occ}}$ contain occupied coefficients with
$C_{\mathrm{occ}}^\dagger S C_{\mathrm{occ}}=I$. The mixed-index
projector and position operators are
\begin{equation}
 \Pi=C_{\mathrm{occ}}C_{\mathrm{occ}}^\dagger S,\qquad
 \mathsf X_a=S^{-1}X_a,\qquad \Pi^2=\Pi.
 \label{eq:mixed-fermi-projector}
\end{equation}
Here $X_a$ denotes the covariant AO matrix of the position operator.
For a physical bulk window $w(\mathbf r)$, define its covariant
matrix $W_{\mu\nu}=\langle\phi_\mu|w|\phi_\nu\rangle$,
$\mathsf W=S^{-1}W$, and the area
$A_w=\int w(\mathbf r)\,\mathrm d^2r$. A finite-window version of
Eq.~\eqref{eq:nc-chern} is
\begin{equation}
 C_w=-\frac{2\pi\ii}{A_w}\,
 \Tr\!\left(\mathsf W\Pi
   [[\mathsf X_1,\Pi],[\mathsf X_2,\Pi]]\right).
 \label{eq:finite-window-chern}
\end{equation}
All mixed-index matrices transform by similarity under a nonsingular
AO basis change. This preserves the trace without assigning a local
contribution to an arbitrary subset of basis coefficients.
The area is a spatial measure, not a count of occupied states or AOs.
In an open finite model with commuting position matrices the trace
over the entire sample cancels between bulk and boundary contributions.
A bulk window instead permits comparison with the local Chern marker
\cite{BiancoResta2011}. A finite value need not be exactly quantized.
Projected AO positions need not commute, so the lattice cancellation
identity must not be imposed on an incomplete continuum basis.

The complex finite-window evaluator uses metric solves and checks
occupied-state orthonormality and $0\leq W\leq S$. Its finite-system
Quickstep interface integrates two-dimensional Gaussian windows
analytically over the AOs, with widths $\sigma_1,\sigma_2$ and area
$A_w=2\pi\sigma_1\sigma_2$. The projector is defined by an explicit
energy in a resolved gap, not by smeared SCF occupations.
Open-boundary model and GPW/GAPW controls are provided in the
Supporting Information. This spatial diagnostic complements rather
than duplicates Wilson-loop or $\mathbb Z_2$ analysis.

Periodic boundaries require a separate treatment. Finite-torus
approximations periodize the spatial derivations rather than applying
an unbounded Cartesian position to a periodic cell
\cite{Prodan2013,BauMarrazzo2024}. Such approximate derivations do not
generally obey an exact Leibniz rule at finite size. Differentiating
the finite Fermi matrix and replacing that derivative by a spectral
Fermi divided difference times a Hamiltonian derivative are therefore
distinct finite-size procedures. Hall and Chern crosschecks must
retain this distinction, the order of tensor indices, and the relative
orientation of the Berry and Clifford constructions.

A spectral gap excludes eigenvalues in an interval. A mobility gap
instead concerns localization and transport even when states are
present. A localizer gap is the smallest singular value of a composite
energy-position operator. These three gaps answer different questions.
The localizer is useful precisely because spatially resolved
information can remain available where a bulk band classification is
not directly applicable \cite{Dixon2023}. Its numerical invertibility
does not by itself prove a mobility-gap theorem for a DFT material.

\subsection{Thermal disorder and convergence}
An AIMD trajectory can supply configurations for a disorder average.
The nuclear sampling temperature, the electronic Fermi temperature,
and the dissipative relaxation scale are conceptually distinct.
A calculation should record all three rather than infer them from a
single input temperature. Configuration correlations determine how
many effectively independent samples have been obtained; a small
standard error computed from strongly correlated snapshots is not a
substitute for a correlation analysis.

The thermodynamic limit, zero-temperature limit, and vanishing
dissipation limit cannot generally be interchanged. At fixed finite
volume the spectrum is discrete, so taking the broadening to zero
first can give a qualitatively misleading DC response. For a material
prediction we therefore require size and sampling convergence at a
specified positive dissipation, followed by a physically justified
study of its dependence. The present validation examples establish
numerical identities at fixed inputs, not a microscopic scattering
model or a disorder-converged transport prediction.
